% =====================================================================
% Section I -- Introduction
% SpecAsync-UVM manuscript, draft 3 (rewrite, not a revision of draft 2)
%
% WHY A REWRITE: draft 2's central claim ("no benefit even under oracle-
% perfect prediction") rested on a misaligned oracle (E0.5, AC-13). The
% pre-registered Gate E series replaced it, and two falsification triggers
% fired (E4: speculation beats the shipped driver; E6: threshold tuning
% matches or beats speculation). E7 confirmed the threshold result; E8
% showed it reverses for sparse oversubscribed access.
%
% REMOVED from draft 2 (superseded; see SUPERSEDED_VALUES.md):
%   - oracle-null lead evidence (p=0.78; +6.7%; +3.60%)      -> CS2-7
%   - spec_hits hit rates (0.0132%, 0.2357%, 99.61%)         -> CS2-N7
%   - dispatch-race explanation (1,685-1,894x; 432-438x)     -> R1
%   - "any such scheme" pipelining ceiling (0.08-7.94%)      -> CS2-15
%   - metadata 0.04 us premise (source file never located)
%   - "<400 lines", "preserving the driver's safety invariants" -> CS2-17
%
% EVIDENCE for every number below (ids in EVIDENCE_EXTRACT.md after
% C3-APPLY): E4 F1/F4, E5 D(t), E6 Family 1 + secondary, E7 Families A/B,
% E8 primary. Platform for all Gate E numbers: RTX 5070 Ti, driver
% 595.91.07 (kernel 7.0.0-31 for E0.5-E0.9, 7.0.0-34 from E0.9a-2 on).
%
% LIT CHECKED 2026-10-10 against the full papers:
%   go2023earlyadaptor: RTX 3090, open driver 515.65.01, CUDA 11.7; thresholds
%     1-100%; in memory (Fig. 4a) 1% >= 51% for every workload; at 150%
%     oversubscription (Fig. 4b) low thresholds hurt irregular workloads
%     (ATA, BIC, NW), explained by thrashing (Sec. III-C/D). Pages 248-258.
%   shen2026stride: user-space LD_PRELOAD layer, CUPTI fault events,
%     cudaMemAdvise + cudaMemPrefetchAsync; A100/V100, CUDA 12.8; up to 1.48x
%     measured vs a "no-prefetch baseline" (driver prefetcher state not
%     stated). ISMM '26 pp. 52-65.
% [T4] filled from GATE_E7T4_REPORT.md (dense-access threshold result only).
%
% OPEN before submission:
%   [SRC]  verify from uvm_migrate.c that cudaMemPrefetchAsync to a GPU
%          installs mappings (MAKE_RESIDENT_AND_MAP), before Section II
%          relies on it.
%   [CAT]  artifact count (25) to be confirmed against ARTIFACT_CATALOG.md.
%   [SEC]  section numbers in the roadmap follow the planned outline.
% =====================================================================

\section{Introduction}

Unified Virtual Memory (UVM) lets GPU kernels address host and device memory
through one pointer space, migrating pages on demand when the GPU touches data
it does not hold. Each such access raises a fault that the host driver must
service: it locates the faulting range, decides where the data should reside,
migrates it, and installs a GPU mapping, largely while holding a lock on the
address space. On fault-intensive workloads this host-side path is a large
share of execution time. NVIDIA's driver already attacks it with an in-path
prefetcher: when it services a fault, it also migrates neighbouring pages if
enough of their surrounding region is already resident, using a density
threshold over a tree of subregions within each 2\,MB block
\cite{ganguly2019interplay}. A natural next step, proposed repeatedly in the
literature and suggested explicitly by Allen \emph{et al.}\ after their
fine-grained decomposition of the UVM fault path \cite{allen2024finegrain}, is
to go further asynchronously: predict which pages will fault next and stage
them \emph{off} the fault path, so the cost is paid before the fault arrives.

This paper tests that idea on real hardware. We built it as
\emph{SpecAsync-UVM}, a modification of NVIDIA's open-source UVM driver in which
a background worker stages pages that a predictor nominates while the
fault-servicing thread continues. To ask what off-path staging can achieve at
all, rather than how well some heuristic predicts, we drive it with a
\emph{first-touch oracle}: a table of the run's own page order recorded in an
earlier run, which on our stencil workload is nearly deterministic, looked up in
constant time. We then measured it against the driver as shipped, under
pre-registered analysis plans with falsification triggers, on an RTX 5070 Ti
(Blackwell) with driver 595.91.07.

The answer has four parts, and the last two concern the driver's own
prefetcher rather than speculation.

\paragraph{Off-path staging cannot prevent a fault}
The background worker can make a page resident, but in this driver only the
fault-servicing path installs GPU mappings. A staged page therefore still
faults when the GPU first touches it. We confirmed this from source, from
counters, and with the oracle: with the prefetcher disabled, staging 99.95\% of
pages ahead of demand left the demand-fault count unchanged, and that
configuration remained $2.92\times$ slower than the shipped driver. Speculation
can move copy work off the fault-servicing thread, worth 21--23\% against a
driver whose prefetcher is off, but it cannot recover what turning the
prefetcher off costs.

\paragraph{Against the shipped driver, speculation helps only through the prefetcher}
With the prefetcher on, whole-block staging with the oracle beats the shipped
driver on a stencil workload by 3.4--5.3\% across four independent sessions,
with no detectable gain on graph traversal. The gain does not come from the
staged pages themselves. The prefetcher's density rule counts \emph{resident}
pages, so pages staged by the worker push regions over its threshold, and the
prefetcher then maps whole regions during ordinary fault service, preventing
about 40\% of demand faults. A pre-registered dose--response test confirms
this: the fault reduction grows from 40\% to 53\% as the threshold is raised from
51 to 75, and vanishes when the rule is disabled.

\paragraph{The same lever is available without speculation}
If speculation's benefit runs through the density threshold, the threshold can
be moved directly. Lowering it from the shipped 51 to 0, with no speculation,
makes the stock driver 9.57\% faster on the same stencil workload (confirmatory
test), and the oracle-driven speculation configuration at the shipped threshold
is 6.4\% slower than it. At matched thresholds, speculation adds at most 2.8\%,
and nothing detectable at threshold 0 or 10 (minimum detectable effects 1.8\%
and 2.0\%). Across nine dense or streaming workloads, the lower threshold sped up
four, by 2.5--17\%, and significantly slowed none.

\paragraph{But the lever cuts both ways}
The shipped threshold is a compromise, not an oversight. On a synthetic
benchmark that touches only one or eight pages per 2\,MB block, the low
threshold is 41--85\% \emph{slower} once the array exceeds GPU memory. Tracing
shows why: at threshold 0 the prefetcher brings in whole regions for a few
useful pages, and on the oversubscribed array a traced run migrates about
130\,GiB into the GPU, against 4\,GiB at the shipped threshold. With the array
in memory the effects are small and mixed, and with 64 or more pages touched per
block the low threshold is faster again. No fixed lower threshold dominates the
shipped one.

\paragraph{Relation to prior work on the threshold}
Neither direction of this effect is new. Go \emph{et al.} varied the same
threshold on an RTX 3090 with release 515.65.01 of the open driver. With memory
not oversubscribed, a low threshold was at least as fast as the default for
every workload they tested; at 150\% oversubscription, the ratio our sparse
test also uses, it hurt irregular workloads through thrashing. They proposed an
adaptive threshold on that basis \cite{go2023earlyadaptor}. Our threshold
measurements replicate both directions on later GPU generations and a later
driver, under pre-registered tests, with access density isolated by a synthetic
benchmark and the migrated bytes measured. What is new is the link to
speculation: the measured value of off-path staging is a side effect of the same
density rule.

\medskip
The result reverses this study's own earlier conclusion, and how it was
overturned is part of what we report. An earlier version concluded that even
oracle prediction yields no benefit. That oracle, we later found, replayed its
trace at the wrong rate: faults were recorded once per batch but consumed once
per fault, so the replay ran ahead of the application and silently wrapped. The
corrected oracle produced the positive result above, and a falsification
trigger committed before the run is what forced us to report it. The
threshold experiment that then explained the positive result away was itself
pre-registered with a trigger, and it fired. Each reversal was caught by a check
written before the measurement it tested.

\subsection{Contributions}

\begin{itemize}
  \item \textbf{An in-driver off-path staging prototype with a usable oracle.}
  A modification of NVIDIA's open-source UVM driver with a constant-time,
  lock-free first-touch oracle and a configurable staging width, whose kernel
  changes passed a review protocol of compile-only review, a precondition table
  for every driver call, and an attended first load.

  \item \textbf{A mechanistic account of what off-path staging can and cannot
  do.} Staging cannot prevent faults because it cannot map (verified from
  source and by counters). Its only route to fault prevention is the
  prefetcher's density rule, confirmed by a pre-registered dose--response
  test. Its other effect, copy offload, is large only when the prefetcher is
  off. We also measure the cost of handing work to the worker
  (54--75\,ns per enqueue on the stencil workload, on the critical path) and show that most of the
  per-fault overhead in our first oracle was its own lookup.

  \item \textbf{An independent replication of the threshold trade-off.}
  Pre-registered tests on the stock driver reproduce Go \emph{et al.}'s finding
  in both directions on a newer driver: the gain from a lower threshold on
  dense and streaming access, on two GPU generations, and the loss on sparse
  access to an oversubscribed array, with the migrated bytes traced and minimum
  detectable effects reported for every null.

  \item \textbf{A falsifiable measurement protocol, exercised.} Every gate fixed
  its families, statistic, test, minimum detectable effect, and falsification
  trigger in a committed file before running. Two triggers fired and changed the
  paper's conclusion. Analysis uses medians, two-sided Mann--Whitney $U$, Holm
  correction across pre-declared families, and reports outliers with and
  without removal.

  \item \textbf{A catalog of UVM measurement artifacts.} Twenty-five artifacts
  met in this study, each with the false conclusion it would have produced,
  including the misaligned oracle. One pattern stands out: errors arising from
  a mismatch of aggregation level (per batch versus per fault, per trial versus
  per run) were caught only by later, independent passes, after they had spread
  into other reports, whereas pre-registered predictions caught errors inside the
  gate that made them.
\end{itemize}

\subsection{Scope}

Our speculation results use a near-perfect oracle on the stencil workload and a
weaker one on graph traversal, so they bound what off-path staging can achieve
and do not evaluate a practical predictor. The dense-access threshold gain
replicates on a Tesla T4 (Turing) with the same driver: $-5.1\%$ on the stencil
workload, and 12 of 14 further tests faster with none slower. The speculation,
sparse-access and oversubscription results are from the RTX 5070 Ti alone. The
sparse-access counterexample uses one synthetic benchmark at one
oversubscription ratio. Prefetching issued from user space is a different
mechanism, and we make no claim about it: Shen and Nikolopoulos predict faults
from profiler events and issue explicit prefetch calls, which the driver
services on its own migration path rather than through an off-path worker,
reporting speedups of up to $1.48\times$ on A100 and V100 GPUs
\cite{shen2026stride}. Section~II discusses the difference.

Section~II reviews related work. Section~III describes the UVM fault path and
the prefetcher's density rule. Section~IV presents SpecAsync-UVM and the oracle,
and Section~V the measurement protocol. Section~VI reports the speculation
results and their mechanism, Section~VII the threshold measurements, and
Section~VIII the artifact catalog. Section~IX concludes.
% [SEC] renumber if the outline changes.
